\documentclass[11pt]{article}
\input{style-sujet}
\usepackage{array}

\begin{document}

\entetesujet{Sujet bac 2018 -- Série D}

% ====================== EXERCICE 1 ======================
\exo{1}{5 points}

Le plan complexe $\Cb$ étant rapporté au repère orthonormal direct $(O,\vec{u},\vec{v})$. On considère les points $A$, $B$, $C$ et $D$ d'affixes respectives
\[ Z_A = 1 + 2i \;;\quad Z_B = -1 + 2i \;;\quad Z_C = 1 - i \;;\quad Z_D = 1 \]

\begin{enumerate}
  \item \begin{enumerate}[label=\alph*.]
    \item Déterminer l'affixe $Z_{\vv{BC}}$ du vecteur $\vv{BC}$.
    \item Déterminer l'expression analytique de la translation de vecteur $\vv{BC}$.
    \item Trouver l'affixe du point $A'$ image du point $A$ par la translation de vecteur $\vv{BC}$.
  \end{enumerate}
  \item \begin{enumerate}[label=\alph*.]
    \item Prouver qu'une mesure, en radian, de l'angle $(\vv{AD},\vv{AB})$ est $-\dfrac{\pi}{2}$.
    \item Écrire l'expression analytique de la rotation $R$ de centre $A$ et d'angle $(\vv{AD},\vv{AB})$.
    \item Trouver l'affixe du point $C'$ image du point $C$ par la rotation $R$.
  \end{enumerate}
  \item Déterminer le rapport et l'angle de la similitude plane directe $S$ de centre $A$ et qui transforme $B$ en $A$.
\end{enumerate}

% ====================== EXERCICE 2 ======================
\exo{2}{5 points}

Soit $\mathcal{E}$ un plan vectoriel rapporté à sa base canonique $(\vi,\vj)$. On considère les deux droites vectorielles $(D_1)$ et $(D_2)$ d'équations cartésiennes respectives $x - 2y = 0$ et $x + y = 0$ de ce plan.

\begin{enumerate}
  \item Vérifier que les droites $(D_1)$ et $(D_2)$ sont engendrées respectivement par les vecteurs $\vec{e_1} = 2\vi + \vj$ et $\vec{e_2} = -\vi + \vj$.
  \item Prouver que la famille $(\vec{e_1},\vec{e_2})$ est une base de $\mathcal{E}$.
  \item Montrer que les sous-espaces vectoriels $(D_1)$ et $(D_2)$ sont supplémentaires dans $\mathcal{E}$.
  \item Soit $f$ un endomorphisme de $\mathcal{E}$ défini par : $f(\vec{e_1}) = \vec{e_1}$ et $f(\vec{e_2}) = -\vec{e_2}$. Exprimer les vecteurs $f(\vi)$ et $f(\vj)$ dans la base $(\vi,\vj)$.
\end{enumerate}

% ====================== EXERCICE 3 ======================
\exo{3}{6 points}

Soit la fonction numérique $f$ à variable réelle $x$, définie par :
\[ f(x) = \begin{cases} \dfrac{\ln x}{-1 + \ln x} & \text{si } x > 0 \\[1.5ex] e^{-2x} & \text{si } x \le 0 \end{cases} \]
On désigne par $(C)$ la courbe représentative de $f$ dans le repère orthonormé $(O,\vi,\vj)$ d'unité graphique : 2 cm.

\begin{enumerate}
  \item Déterminer l'ensemble de définition de $f$.
  \item Vérifier que la fonction $f$ est continue en $x = 0$.
  \item Étudier la dérivabilité de $f$ en $x = 0$.
  \item \begin{enumerate}[label=\alph*.]
    \item Déterminer la fonction dérivée $f'$ de $f$.
    \item Dresser le tableau de variation de $f$.
  \end{enumerate}
  \item \begin{enumerate}[label=\alph*.]
    \item Préciser les branches infinies à la courbe $(C)$ de $f$.
    \item Tracer $(C)$.
  \end{enumerate}
  \item Calculer l'aire $\mathcal{A}$ du domaine limité par la courbe $(C)$ de $f$, l'axe des abscisses et les droites d'équations $x = -1$ ; $x = 0$.
\end{enumerate}

% ====================== EXERCICE 4 ======================
\exo{4}{4 points}

Le tableau ci-dessous représente le couple $(x, y)$ des deux caractères d'une série statistique. $x$ est le nombre de jours et $y$ le poids en mg d'une larve.

\begin{center}
\renewcommand{\arraystretch}{1.3}
\begin{tabular}{|c|c|c|c|c|c|c|}
\hline
$x$ & 1 & 2 & 3 & 4 & 5 & 6 \\
\hline
$y$ & 0,2 & 1,4 & 1,8 & 2 & 2,6 & 3 \\
\hline
\end{tabular}
\end{center}

\begin{enumerate}
  \item Calculer les coordonnées $\overline{x}$ et $\overline{y}$ du point moyen $G$.
  \item Déterminer l'équation de la droite de régression linéaire de $y$ en $x$.
  \item Estimer le poids de la larve au 7\textsuperscript{e} jour.
\end{enumerate}

\end{document}
